\documentclass[11pt,a4paper]{article}
\usepackage[utf8]{inputenc}
\usepackage[T1]{fontenc}
\usepackage[ngerman]{babel}
\usepackage{helvet}
\renewcommand{\familydefault}{\sfdefault}
\usepackage[left=2cm, right=2cm, top=2cm, bottom=2cm]{geometry}
\usepackage{amsmath,amssymb}
\usepackage{array}
\usepackage{booktabs}
\usepackage{tikz}
\usepackage{pgfplots}
\pgfplotsset{compat=1.18}
\usepackage{fancyhdr}
\usepackage{xcolor}

% Minimal Farben
\definecolor{border}{RGB}{221,221,221}
\definecolor{accent}{RGB}{44,90,160}
\definecolor{correct}{RGB}{34,139,94}

\pagestyle{fancy}
\fancyhf{}
\fancyhead[L]{\small Physik Klasse 10E}
\fancyhead[R]{\small\color{correct}\bfseries ERWARTUNGSHORIZONT}
\fancyfoot[C]{\small Seite \thepage}
\renewcommand{\headrulewidth}{0.4pt}

\setlength{\parindent}{0pt}
\newcounter{aufgabe}
\newcommand{\aufgabe}[1]{%
\stepcounter{aufgabe}%
\vspace{0.5cm}%
{\bfseries Aufgabe \theaufgabe}\hfill{\small #1 BE}\\[0.2cm]%
}

\newcommand{\lsg}[1]{\textcolor{correct}{#1}}
\newcommand{\be}[1]{\hfill\textcolor{correct}{\small(#1 BE)}}

\begin{document}

% === KOPFBEREICH ===
\begin{center}
{\Large\bfseries Probeklausur: Kinematik}\\[0.2cm]
{\large\color{correct}\bfseries Erwartungshorizont}\\[0.3cm]
{\small Klasse 10E \quad|\quad 45 Minuten \quad|\quad 36 BE}
\end{center}

\vspace{0.3cm}
{\color{border}\hrule height 0.5pt}
\vspace{0.3cm}

% === NOTENSPIEGEL ===
\begin{center}
\small
\begin{tabular}{|c|c|c|c|c|c|c|}
\hline
\textbf{Note} & 1 & 2 & 3 & 4 & 5 & 6 \\
\hline
\textbf{ab BE} & 35 & 29 & 22 & 15 & 7 & 0 \\
\hline
\textbf{ab \%} & 96 & 80 & 60 & 40 & 20 & 0 \\
\hline
\end{tabular}
\end{center}

\vspace{0.2cm}
{\color{border}\hrule height 0.5pt}
\vspace{0.5cm}

% === LÖSUNGEN ===

\aufgabe{6}
\textbf{Grundlagen} (AFB I)

\begin{enumerate}
\item[a)] \lsg{90 km/h = 90 $\div$ 3,6 = 25 m/s} \be{1}

\item[b)] \lsg{Formel: $s = 5t^2$} \be{1}\\
\lsg{Einsetzen: $s = 5 \cdot 4^2 = 5 \cdot 16 = 80\,\mathrm{m}$} \be{1}

\item[c)] Tabelle: \be{3}

\begin{center}
\begin{tabular}{|l|c|c|}
\hline
\textbf{Diagramm} & \textbf{Steigung bedeutet} & \textbf{Fläche bedeutet} \\
\hline
$s(t)$-Diagramm & \lsg{Geschwindigkeit $v$} & --- \\
\hline
$v(t)$-Diagramm & \lsg{Beschleunigung $a$} & \lsg{Weg $s$} \\
\hline
\end{tabular}
\end{center}
{\small\color{correct} Je 1 BE pro richtige Eintragung}
\end{enumerate}

\vspace{0.5cm}

\aufgabe{8}
\textbf{Bremsvorgang} (AFB II)

\begin{enumerate}
\item[a)] \lsg{$v_0 = 72\,\mathrm{km/h} = 72 \div 3{,}6 = 20\,\mathrm{m/s}$} \be{1}

\item[b)] \lsg{Formel: $t = \dfrac{v_0}{|a|}$} \be{1}\\[0.2cm]
\lsg{Einsetzen: $t = \dfrac{20\,\mathrm{m/s}}{4\,\mathrm{m/s^2}} = 5\,\mathrm{s}$} \be{1}

\item[c)] \lsg{Formel: $s = \dfrac{v_0^2}{2|a|}$} \be{1}\\[0.2cm]
\lsg{Einsetzen: $s = \dfrac{(20\,\mathrm{m/s})^2}{2 \cdot 4\,\mathrm{m/s^2}} = \dfrac{400}{8}\,\mathrm{m} = 50\,\mathrm{m}$} \be{2}

\item[d)] \lsg{Ja, der Bremsweg (50 m) ist kürzer als der Abstand zum Hindernis (60 m).} \be{1}\\
\lsg{Das Auto hält 10 m vor dem Hindernis.} \be{1}
\end{enumerate}

\vspace{0.5cm}

\aufgabe{8}
\textbf{Freier Fall und senkrechter Wurf} (AFB II)

\begin{enumerate}
\item[a)] \lsg{Formel: $t_{\text{steig}} = \dfrac{v_0}{g}$} \be{1}\\[0.2cm]
\lsg{Einsetzen: $t_{\text{steig}} = \dfrac{20\,\mathrm{m/s}}{10\,\mathrm{m/s^2}} = 2\,\mathrm{s}$} \be{1}

\item[b)] \lsg{Formel: $h_{\text{max}} = \dfrac{v_0^2}{2g}$} \be{1}\\[0.2cm]
\lsg{Einsetzen: $h_{\text{max}} = \dfrac{(20\,\mathrm{m/s})^2}{2 \cdot 10\,\mathrm{m/s^2}} = \dfrac{400}{20}\,\mathrm{m} = 20\,\mathrm{m}$} \be{2}

\item[c)] \lsg{Der Ball kommt mit $v = 20\,\mathrm{m/s}$ (nach unten) an.} \be{1}\\
\lsg{Begründung: Die Energie bleibt erhalten. / Die zeitlose Formel liefert am Startpunkt wieder $|v| = v_0$. / Die Bewegung ist symmetrisch.} \be{2}

{\small\color{correct} 2 BE für korrekte Begründung, 1 BE für richtige Geschwindigkeit}
\end{enumerate}

\newpage

\aufgabe{6}
\textbf{Diagramm-Interpretation} (AFB II)

\begin{enumerate}
\item[a)] \lsg{Phase 1 (0-4 s): Beschleunigte Bewegung (Geschwindigkeit nimmt zu)} \be{1}\\
\lsg{Phase 2 (4-10 s): Gleichförmige Bewegung (Geschwindigkeit konstant 8 m/s)} \be{1}

\item[b)] \lsg{Formel: $a = \dfrac{\Delta v}{\Delta t}$} \be{1}\\[0.2cm]
\lsg{Einsetzen: $a = \dfrac{8\,\mathrm{m/s} - 0}{4\,\mathrm{s}} = 2\,\mathrm{m/s^2}$} \be{1}

\item[c)] \lsg{Fläche = Dreieck + Rechteck} \be{1}\\
\lsg{$s = \dfrac{1}{2} \cdot 4 \cdot 8 + 6 \cdot 8 = 16 + 48 = 64\,\mathrm{m}$} \be{1}

{\small\color{correct} Alternativ korrekt: $s_1 = \frac{1}{2}at^2 = \frac{1}{2} \cdot 2 \cdot 16 = 16\,\mathrm{m}$ und $s_2 = v \cdot t = 8 \cdot 6 = 48\,\mathrm{m}$}
\end{enumerate}

\vspace{0.5cm}

\aufgabe{8}
\textbf{Luftwiderstand und Transfer} (AFB III)

\begin{enumerate}
\item[a)] \lsg{Beim Öffnen des Fallschirms vergrößert sich die Querschnittsfläche $A$ stark.} \be{1}\\
\lsg{Da $F_W \sim A$ (und $F_W \sim v^2$), wird der Luftwiderstand viel größer.} \be{1}\\
\lsg{Der Fallschirmspringer erreicht eine niedrigere Grenzgeschwindigkeit.} \be{1}

{\small\color{correct} Relevante Größen: Fläche $A$, Luftwiderstandskraft $F_W$, Grenzgeschwindigkeit $v_{\text{grenz}}$}

\item[b)] \lsg{In 39 km Höhe ist die Luftdichte $\rho$ viel geringer als am Boden.} \be{1}\\
\lsg{Weniger Luftdichte bedeutet weniger Luftwiderstand.} \be{1}\\
\lsg{Dadurch ist die Grenzgeschwindigkeit höher.} \be{1}

{\small\color{correct} Akzeptabel auch: „fast Vakuum", „kaum Luft", „Luft 300× dünner"}

\item[c)] \lsg{Die Aussage ist nur teilweise richtig:} \be{0}\\
\lsg{\textbf{Ohne Luftwiderstand (Vakuum):} Falsch. Alle Körper fallen gleich schnell, unabhängig von der Masse (Galilei). $s = \frac{1}{2}gt^2$ enthält keine Masse.} \be{1}\\
\lsg{\textbf{Mit Luftwiderstand:} Tendenziell richtig. Schwerere Körper haben oft eine höhere Grenzgeschwindigkeit, da sie mehr Masse pro Fläche haben.} \be{1}

{\small\color{correct} Für 2 BE: Unterscheidung mit/ohne Luftwiderstand + je eine korrekte Begründung}
\end{enumerate}

\vspace{0.5cm}
{\color{border}\hrule height 0.5pt}
\vspace{0.3cm}

\begin{center}
\textbf{Zusammenfassung der Anforderungsbereiche}\\[0.3cm]
\begin{tabular}{|l|c|c|c|}
\hline
\textbf{AFB} & \textbf{Aufgaben} & \textbf{BE} & \textbf{Anteil} \\
\hline
I (Reproduktion) & 1 & 6 & 17\% \\
\hline
II (Anwendung) & 2, 3, 4 & 22 & 61\% \\
\hline
III (Transfer) & 5 & 8 & 22\% \\
\hline
\textbf{Gesamt} & & \textbf{36} & 100\% \\
\hline
\end{tabular}
\end{center}

\end{document}
